% Responsável: P1 (base legal verificada em 17/11/2026; rever a redação)
\section{Considerações éticas e limitações}\label{sec:etica}

\paragraph{Dados pessoais (RGPD).}
As notícias identificam pessoas reais, sobretudo figuras públicas, e associam-nas a partidos e
posições políticas. Dados que revelem opiniões políticas são uma categoria especial cujo tratamento
o RGPD proíbe em regra (art.~9.º, n.º~1)~\cite{rgpd2016}. Duas exceções do n.º~2 cobrem o nosso uso:
a alínea~e), para dados manifestamente tornados públicos pelo titular, como as declarações públicas
de políticos, e a alínea~j), para fins de investigação científica, com as garantias do art.~89.º.
Algumas notícias mencionam também pessoas privadas (vítimas, familiares, uma criança), e várias
notícias falsas fazem acusações difamatórias. Medidas adotadas: não redistribuímos os textos (o
código descarrega-os da fonte); não construímos perfis de pessoas; no artigo citamos apenas excertos
curtos e, quando possível, descrevemos os casos sem nomes. Os modelos guardam vocabulário, incluindo
nomes próprios, pelo que um sistema em produção deveria minimizar esses atributos e fazer uma
avaliação de impacto sobre a proteção de dados antes de ser usado.

\paragraph{Licença.}
O repositório do corpus não tem ficheiro de licença; os autores pedem apenas a citação dos artigos
de origem. Usamo-lo para investigação académica, sem redistribuir os dados.

\paragraph{Enviesamentos.}
(1)~\emph{Político}: 58\% do corpus é política, e o tópico mais associado às notícias falsas é dominado
por ataques a figuras de um só partido (\emph{Críticas a Lula e ao PT}, T10). Um modelo pode aprender
a associar o vocabulário de um quadrante político à falsidade e penalizar notícias verdadeiras sobre
esses temas, como sugerem os falsos positivos da categoria \emph{Tema típico da outra classe}.
(2)~\emph{De fonte}: 93\% das falsas vêm de um único sítio e quase todas as verdadeiras de três grupos
de comunicação; parte do ``estilo'' aprendido pode ser o estilo editorial desses sítios, como
Potthast et~al.~\cite{potthast2018} alertam. A análise de erros mostra o reverso: notícias reais
copiadas por sítios de desinformação ficam rotuladas como falsas. (3)~\emph{Temporal e linguístico}:
as notícias são de 2016--2018 e em português do Brasil; o vocabulário político muda, e o desempenho
em notícias portuguesas não foi avaliado.

\paragraph{Riscos de utilização indevida.}
Usado como juiz automático, e não como apoio à triagem, o sistema poderia levar à remoção de conteúdo
legítimo (falsos positivos), por exemplo artigos de opinião verdadeiros, que são o erro mais comum do
BERTimbau, ou dar uma falsa sensação de segurança (falsos negativos). Quem produz desinformação pode
imitar o estilo jornalístico, incluindo datas e fontes, para escapar à deteção: é exatamente o que
acontece nos 12 falsos negativos analisados. Por isso, a decisão final deve caber sempre a um
verificador humano.

\paragraph{Nível de risco no Regulamento da IA.}
O Regulamento (UE) 2024/1689~\cite{aiact2024}, alterado pelo Regulamento (UE)
2026/1744~\cite{omnibus2026}, classifica como sistemas de risco elevado os usos do Anexo~III:
biometria, infraestruturas críticas, educação, emprego, serviços essenciais, aplicação da lei,
migração, justiça e processos democráticos. O Anexo~III não inclui a verificação de factos, a
moderação de conteúdos nem a deteção de desinformação; o ponto mais próximo, 8.b), refere sistemas
destinados a influenciar o resultado de eleições ou o comportamento dos eleitores, o que não é o caso
de uma ferramenta interna que ordena notícias para verificação humana. O sistema também não gera
conteúdos, pelo que as obrigações de transparência do art.~50.º não se aplicam. Segundo a nossa
leitura, trata-se portanto de um sistema de risco mínimo. A classificação deveria ser revista se o
sistema passasse a decidir sozinho a remoção de conteúdos ou fosse usado em contexto eleitoral; nesse
caso, a remoção por plataformas em linha está também sujeita ao Regulamento dos Serviços
Digitais~\cite{dsa2022}. Esta análise não substitui um parecer jurídico.

\paragraph{Limitações.}
O corpus é de um só país, período e conjunto de fontes; tem textos quase vazios e alguns rótulos
duvidosos; a regra de escolha de $K$ foi decidida após ver a curva de coerência; os atributos de
estilo do corpus foram calculados sobre o texto completo; a versão cortada reduz algumas notícias a
um título; e a avaliação da extração de informação assenta em 100 documentos sem adjudicação entre
anotadores.
